\documentclass[12pt]{article}
\usepackage{amsmath,amssymb,amsfonts}
\usepackage[margin=1in]{geometry}

\begin{document}

\section*{Chapter 3, Problem 19}

\textbf{Problem.} Prove that the determinant of any matrix is equal to the product of its eigenvalues and that the trace of a matrix is equal to the sum of its eigenvalues. Thus show that if a matrix has a zero eigenvalue it is not invertible.

\bigskip
\textbf{Solution.}

Let $A$ be an $n\times n$ matrix with eigenvalues $\lambda_1, \lambda_2, \ldots, \lambda_n$ (counted with algebraic multiplicity). The characteristic polynomial is:
\[
\det(A - \lambda I) = (-1)^n(\lambda - \lambda_1)(\lambda - \lambda_2)\cdots(\lambda - \lambda_n).
\]

\textbf{Determinant equals product of eigenvalues:} Set $\lambda = 0$:
\[
\det(A) = \det(A - 0 \cdot I) = (-1)^n(-\lambda_1)(-\lambda_2)\cdots(-\lambda_n) = (-1)^n(-1)^n\lambda_1\lambda_2\cdots\lambda_n = \lambda_1\lambda_2\cdots\lambda_n.
\]

\textbf{Trace equals sum of eigenvalues:} Expand the characteristic polynomial. The coefficient of $(-\lambda)^{n-1}$ in $\det(A - \lambda I)$ can be obtained two ways:

From the product form:
\[
(-1)^n(\lambda - \lambda_1)\cdots(\lambda - \lambda_n) = (-1)^n\left[\lambda^n - (\lambda_1 + \cdots + \lambda_n)\lambda^{n-1} + \cdots\right].
\]
So the coefficient of $\lambda^{n-1}$ is $(-1)^n \cdot (-\sum_i \lambda_i) = (-1)^{n+1}\sum_i \lambda_i$.

From the determinant expansion, the term $\lambda^{n-1}$ in $\det(A - \lambda I)$ arises only from the product of diagonal elements:
\[
\prod_{i=1}^{n}(A_{ii} - \lambda) = (-\lambda)^n + (-\lambda)^{n-1}\sum_{i=1}^{n} A_{ii} + \cdots
\]
The coefficient of $\lambda^{n-1}$ is $(-1)^{n-1}\sum_i A_{ii} = (-1)^{n-1}\operatorname{tr}(A)$.

Wait, let me be more careful. We have $(-\lambda)^{n-1} = (-1)^{n-1}\lambda^{n-1}$, so:
\[
\prod_i (A_{ii}-\lambda) = (-1)^n\lambda^n + (-1)^{n-1}\left(\sum_i A_{ii}\right)\lambda^{n-1} + \cdots
\]
The coefficient of $\lambda^{n-1}$ from this is $(-1)^{n-1}\operatorname{tr}(A)$. No other terms in the determinant expansion contribute to $\lambda^{n-1}$ (they are at most degree $n-2$).

Equating: $(-1)^{n-1}\operatorname{tr}(A) = (-1)^{n+1}\sum_i \lambda_i = (-1)^{n-1}\sum_i \lambda_i$.

Therefore $\operatorname{tr}(A) = \sum_{i=1}^{n}\lambda_i$.

\medskip
\textbf{Singular if zero eigenvalue:} If any $\lambda_k = 0$, then
\[
\det(A) = \lambda_1 \cdots \lambda_k \cdots \lambda_n = 0,
\]
so $A$ is singular (not invertible). \hfill $\blacksquare$

\end{document}
